% RESEARCH_DRAFT_SHA256: ac77626c2efb82e982a489c36ca3087c4f9303caeab61a231867c9b92836d42c
% Layout and source-reference wrappers expanded for the web renderer.
\documentclass[11pt]{article}
\usepackage[margin=25mm]{geometry}
\usepackage{amsthm}
\usepackage{fontspec,unicode-math}
\setmainfont{Latin Modern Roman}
\setmathfont{Latin Modern Math}
\usepackage{amsmath,amssymb,mathtools,xurl,hyperref}
\hypersetup{colorlinks=true,urlcolor=blue,linkcolor=blue}
\setcounter{secnumdepth}{0}
\setlength{\parindent}{0pt}
\setlength{\parskip}{5pt}
\setlength{\emergencystretch}{3em}
\newcommand{\sref}[2]{\hyperlink{source:#1:#2}{[S#2]}}
\newcommand{\cataloguescope}[1]{\paragraph{Source record.}#1}
\newcommand{\floor}[1]{\left\lfloor#1\right\rfloor}
\newtheorem{proposition}{Proposition}
\newtheorem{lemma}[proposition]{Lemma}
\begin{document}
\paragraph{Research manuscript.} OpenAI \texttt{gpt-6-astra}, reasoning \texttt{ultra}; 11 October 2026.
\paragraph{Preserved source body.} The complete catalogue mathematical text and sources follow with layout and source links expanded, including source qualifications and its historical preparation notes. New research begins at the section headed Research attempt.

\section{Explicit Equations Separating Matrix Multiplication from Low-Rank Tensors}\label{statement}

\subsection{Definitions and mathematical statement}

\paragraph{Shared notation.}Standard number systems are $\mathbb N,\mathbb Z,\mathbb Q,\mathbb R,\mathbb C$. Specialized notation and index ranges follow the cited source and the definitions below.


The rank of a matrix is the dimension of its column space. Matrix completion reconstructs a low-rank matrix from a subset of its entries; matrix sensing reconstructs it from specified linear measurements. Tensor decomposition represents a multiway array as a sum of simple outer products.
\paragraph{Record scope and qualifications.}
The catalogue title is retained as metadata. Named certificates, contracts, bounds or frameworks in that title are not asserted as results of the original source extract.
The supplied registry flags this extraction as unrecoverable. It is displayed as a literal source fragment; the definitions provide context rather than a replacement formal statement.

Source-record limitation: Source flags the statement as unrecoverable: The official report poses this as a rhetorical transition and immediately answers it by introducing representation-theoretic methods for finding equations in the border-rank ideal. No separate clean statement is promoted; the original record remains publicly visible for provenance.
The source text below is rendered literally where necessary; damaged formulas and composite extracts are not silently repaired into a new mathematical claim.
The supplied catalogue marks this record as unrecoverable. Its extract and terminology are documented, but a complete original problem has not been reconstructed.
\paragraph{Statement recorded in the supplied source.}
Given the variety of tensors of border rank at most r, how can one find explicit polynomials in its ideal that vanish on all such tensors but do not vanish on the matrix-multiplication tensor? \href{https://doi.org/10.4171/owr/2018/51}{[S2]}
\paragraph{Source status.}
The supplied catalogue labels this source record open; that metadata was checked against the supplied snapshot on 2026-10-11. The label is preserved for provenance and is not an independent verification of present open status or mathematical correctness. The catalogue literature assessment is dated 2026-08-21.
\subsection{Short English statement}
The supplied fragment concerns explicit Equations Separating Matrix Multiplication from Low-Rank Tensors. Its complete mathematical target is not recovered.
\subsection{Sources}
\begin{itemize}
\item[S1] Ulam AI and the UnsolvedMath Contributors, UnsolvedMath: A Curated Collection of Open Mathematics Problems, catalogue entry. Catalogue curation and metadata: CC BY 4.0. The source extract is retained; definitions, English summaries and qualifications are editorial additions.
\url{https://huggingface.co/datasets/ulamai/UnsolvedMath}.
\item[S2] Complexity Theory (2019), Oberwolfach Reports, DOI 10.4171/owr/2018/51, pp. 3025--3080
\url{https://doi.org/10.4171/owr/2018/51}.
\item[S3] Oberwolfach Report 51/2018
\url{https://ems.press/content/serial-article-files/46773}.
\item[S4] More Asymmetry Yields Faster Matrix Multiplication
\url{https://doi.org/10.1137/1.9781611978322.63}.
\item Source-specific Oberwolfach Reports licences, including CC BY-SA 4.0 where applicable, are retained; see THIRD\_PARTY\_NOTICES.md.
\end{itemize}
\paragraph{Preparation and review.}
The original source and UnsolvedMath contributors retain attribution for the source extract. Codex (GPT-6; exact serving revision not exposed) assisted with definitions, English summaries, scope notes and formatting. No new human mathematical authorship or independently verified solution is asserted. Catalogue identity checked against the complete supplied registry and current repository. Distinct upstream records, including repeated questions and broad or damaged extracts, are retained for complete coverage of the source CS category. A repeated record is not a new independent conjecture. Complete mathematical-equivalence and current-literature review has not been performed. Preparation and final review follow contribution prompts C and D, with source gaps disclosed. Structural and rendering checks do not establish a new theorem or open status.
\label{end:statement}

\section{Research attempt}
\subsection{Literature, source context and scope}
The full original catalogue body and source qualifications are preserved above. The relevant primary workshop passages were read in the cached publisher report. Historical claims and subsequent attributed results are distinguished from the mathematical work below. A restricted example, standard lemma or known theorem is not counted as new progress on the entire original target. Missing assumptions and repeated or rhetorical extracts are not silently turned into different solved problems.

\subsection{An explicit flattening equation reaches only a basic rank threshold}
The primary passage is a rhetorical transition immediately followed by representation-theoretic methods. It does not specify a new residual equation-construction target. We work out a simple instance of the proposed method and locate its limitation. Over $\mathbb C$, write the matrix-multiplication tensor as
\[
T_n=\sum_{i,j,\ell=1}^n e_{ij}\otimes e_{j\ell}\otimes e_{\ell i}.
\]
Flatten a tensor by using its first factor as row index and the pair of remaining factors as column index. This is an $n^2$ by $n^4$ matrix.
\begin{proposition}
The determinant of the $n^2$ by $n^2$ submatrix using columns $(e_{j1},e_{1i})$, indexed by pairs $(i,j)$, vanishes on all tensors of border rank at most $n^2-1$, but equals 1 at $T_n$.
\end{proposition}
\begin{proof}
A simple tensor $a\otimes b\otimes c$ flattens to the rank-at-most-one matrix $a(b\otimes c)^T$. A sum of $r$ simple tensors therefore has flattening rank at most $r$, and every $(r+1)$-minor vanishes. The same holds for border rank at most $r$, since these minors are polynomials, hence continuous, and vanish on every approximating rank-at-most-$r$ tensor.
In $T_n$, the coefficient in row $(i',j')$ and selected column $(e_{j1},e_{1i})$ is 1 exactly when $i'=i$ and $j'=j$, and is otherwise 0. The selected submatrix is therefore the identity when its rows and columns use the same pair order. Its determinant is 1, proving the assertion.
\end{proof}
The failed extension is trying to obtain larger lower bounds from minors of this same flattening. It has only $n^2$ rows, so no minor of order greater than $n^2$ exists. For a proposed border-rank threshold $r\ge n^2$, this particular rank test imposes no further constraint. The source already discusses substantially stronger methods; the displayed determinant is a basic known equation, not an improvement on matrix-multiplication bounds.
\subsection{Verification and remaining gaps}
The coefficient tensor and selected minor were constructed for $n=1,2,3$, producing identity matrices of orders1,4,9. Small sums of simple tensors were also flattened to check the expected minor vanishing. No new high-rank separating polynomial or general representation-theoretic construction is obtained. The original unrecoverable-fragment designation and full source context remain preserved.

\section{References and source attribution}
Primary workshop report: \url{https://doi.org/10.4171/owr/2018/51}. Relevant passages read: 2018/51 pp. 3043--3044, border rank, ideal equations, rhetorical transition and representation-theoretic/minor methods.

All retrievals were checked on 11 October 2026. Original question attribution remains in the preserved Sources section. External theorems are attributed and not claimed as new discoveries or fully verified merely by citation.


\paragraph{Integration review.} Codex (GPT-6; exact serving revision not exposed) reviewed the written arguments and expanded layout and source-reference wrappers. Repository metadata and the production web renderer were checked using Python 3.13 and JavaScript with MathJax 3.2.2. This model-assisted review is not independent human proof verification; no PDF compilation or formal proof checking is asserted.

\section{Completion Estimate}
\noindent\textbf{COMPLETION: 0\%}
No substantiated new progress toward the entire original target is claimed. The worked results and failed route are delimited above; a known result, illustrative model, or auxiliary obstruction is not counted as original-target completion.

\section{Final status and attribution}
\textbf{Attempt status: unresolved.} This describes this research attempt, not a new ruling on the source catalogue's status or on external literature claims.
The mathematics and this self-audit were generated by OpenAI \texttt{gpt-6-astra}, reasoning \texttt{ultra}, on 11 October 2026. Codex (GPT-6; exact serving revision not exposed) orchestrated the batch, reviewed the written mathematical arguments, and contributed formatting/repository review. This additional model review is not independent human verification. No human mathematical authorship, independent human proof checking, or formal verification is claimed. The historical preparation paragraphs in the reproduced catalogue body concern that body only.

Contribution prompt A was applied to the complete definition and supplied source record, followed by the adversarial checks of prompt B. Local reproduction files are \texttt{batch022/checks.py}, \texttt{computations.json}, and the per-ID audit JSON. They were not public links at preparation time. Numerical checks are finite and do not prove unrestricted results. No LaTeX engine was found locally; compilation was not performed by this research agent.

\end{document}
